% English translation of checked_pages/esq_119_009.tex; source: cote 119, site 9, PDF page 10, typescript page 7.
% AI-translated by GPT-6 Astra (OpenAI, ChatGPT Pro mode).
\scanpage{119/9}
\mspage{7}

\unsure{moduli}\tnote{TN-009-01}{The checked French retains the alternatives ``mcdulaires|modulaires'' for the unclear typed letter. Both lead to ``moduli'' in this phrase; the doubt remains marked and linked to the scan.}
multiplicities $M_{g,\nu}$ in the complex analytic context,
namely that, for a moduli dimension $N \geq 3$, the fundamental
group of $M_{g,\nu}$ (i.e. the usual Teichmüller group $T_{g,\nu}$)
is isomorphic to the ``fundamental group at infinity'', i.e. that of a
``tubular neighborhood of infinity''. This is a very familiar fact
(due essentially to Lefschetz) for a smooth \uline{affine} variety of
dimension $N \geq 3$. It is true that the moduli multiplicities are
not affine (except for small values of $g$), but it would suffice
for such an $M_{g,\nu}$ of dimension $N$ (or rather, a suitable finite
covering) to be the union of $N$-2 affine open sets, thus
for $M_{g,\nu}$ not to be ``too close to a compact variety''.

Having no doubt about this principle for constructing the
Teichmüller tower, I prefer to leave it to the experts in the transcendental
theory, better equipped than I am, to prove what is needed
(should any be interested), and instead to make explicit,
with all the care it deserves, the resulting structure of the
Teichmüller tower by generators and relations, this time in the
discrete setting and not the profinite one - which amounts, essentially,
to a complete understanding of the four moduli multiplicities
$M_{0,4}$, $M_{1,1}$, $M_{0,5}$, $M_{1,2}$, and of their fundamental groupoids relative
to suitably chosen ``base points''. These present themselves
quite naturally as the complex algebraic curves of the type
$(g,\nu)$ under consideration whose automorphism group (necessarily
finite) is larger than in the generic case (*). Including among them the
holomorphic sphere with three marked points (coming from $M_{0,3}$, i.e.

\rule{4cm}{0.4pt}

(*) One must also add the ``base points'' arising from operations
of gluing ``pieces'' of the same type in lower moduli dimension.
On the other hand, in moduli dimension 2 (the cases of $M_{0,5}$ and $M_{1,2}$),
one should exclude the points of certain one-parameter families
of curves admitting an exceptional automorphism of order 2. These
families, moreover, constitute remarkable rational curves on the multiplicities
under consideration, which seem to me an important ingredient
of the structure of these multiplicities.

\margin{9}

